\subsection{Local maps to normed spaces}\label{subsec:spectral-local-models}
\paraid{p-160c1dd745e3}
We now turn the spectral subspace into coordinates in
$\RealNumbers^\ModelDimension $.
A choice of orthonormal basis determines the coordinates,
but the resulting pseudometric is independent of this choice.

\paraid{p-643ccba9258d}
We use the orthogonal group
\[
 \OrthogonalGroup(\ModelDimension )=\{\OrthogonalChange \in\RealNumbers^{\ModelDimension \times \ModelDimension }\mid \OrthogonalChange ^{\TransposeSymbol }\OrthogonalChange =\IdentityMatrix _\ModelDimension \},
\]
acting on column vectors in
$\RealNumbers^\ModelDimension $.
Thus
$\OrthogonalChange $
preserves the Euclidean norm and
$\OrthogonalChange ^{-1}=\OrthogonalChange ^{\TransposeSymbol }$.

\paraid{p-f848e928ba37}
We consider the equivalence class of a coordinate map and its norm,
so that a change of orthonormal basis does not change the model.

\begin{definition}\label{def:normed-coordinate-class}
\paraid{p-49a6a71b7246}
Let
$\ComparisonCarrier$
be a nonempty compact space and let
$\ModelDimension\geq1$.
Write
\[
 \FiniteNormSpace{\ModelDimension}
 =\{\NormSymbol\colon\RealNumbers^\ModelDimension\to[0,\infty)
       \mid\NormSymbol\text{ is a norm}\}.
\]
The group
$\OrthogonalGroup(\ModelDimension)$
acts on
$\ContinuousFunctions(\ComparisonCarrier,\RealNumbers^\ModelDimension)
 \times\FiniteNormSpace{\ModelDimension}$
by
\begin{equation}\label{eq:frame-action}
 \OrthogonalChange\cdot(\FeatureCoordinates,\NormSymbol)
 =(\OrthogonalChange\circ\FeatureCoordinates,
   \NormSymbol\circ\OrthogonalChange^{-1}).
\end{equation}
We call an element of the quotient set
\begin{equation}\label{eq:normed-coordinate-class-space}
 \CoordinateClassSpace{\ModelDimension}{\ComparisonCarrier}
 =\bigl(\ContinuousFunctions(\ComparisonCarrier,\RealNumbers^\ModelDimension)
       \times\FiniteNormSpace{\ModelDimension}\bigr)
       /\OrthogonalGroup(\ModelDimension)
\end{equation}
an \emph{$\ModelDimension$-dimensional normed coordinate class on
$\ComparisonCarrier$}.
The class of a pair
$(\FeatureCoordinates,\NormSymbol)$
is denoted by
$[\FeatureCoordinates,\NormSymbol]_{\OrthogonalGroup(\ModelDimension)}$.
\end{definition}

\paraid{p-f230f9fc55e9}
For an integer
$\ModelDimension\geq1$,
a nonempty compact space
$\ComparisonCarrier$,
a continuous map
$\FeatureCoordinates\colon\ComparisonCarrier\to\RealNumbers^\ModelDimension$,
a norm
$\NormSymbol\in\FiniteNormSpace{\ModelDimension}$,
and points
$\ComparisonPoint,\IntegrationPoint\in\ComparisonCarrier$,
set
\[
 \Pseudometric(\ComparisonPoint,\IntegrationPoint)
 =\NormSymbol(\FeatureCoordinates(\ComparisonPoint)-\FeatureCoordinates(\IntegrationPoint)).
\]
This formula defines a continuous pseudometric on
$\ComparisonCarrier$.
The pseudometric is unchanged by \eqref{eq:frame-action}.
Thus it depends only on the normed coordinate class.
Coordinate classes are defined on concrete carriers.
To compare classes on different carriers,
we use the compatibility with input isometries in \Cref{thm:local-models}.
That theorem constructs such a class for each compact metric space
in a neighborhood of a fixed center.
Representing coordinate pairs can be chosen to converge along each convergent sequence of spaces.
We use their linear coordinates to combine approximations by a partition of unity
in Section~\ref{sec:ar}.
Figure~\ref{fig:orthogonal-coordinate-change} illustrates the invariance
under \eqref{eq:frame-action}.

\begin{figure}[H]
\centering
\begin{tikzpicture}[>=Stealth,font=\small]
\node at (0,1.65)
 {$(\FeatureCoordinates,\NormSymbol)$};
\node at (5.8,1.65)
 {$(\OrthogonalChange\circ\FeatureCoordinates,
     \NormSymbol\circ\OrthogonalChange^{-1})$};
\begin{scope}
 \filldraw[fill=black!4] (0,0) ellipse (1.45 and 0.88);
 \draw[->,black!35] (-1.8,0) -- (1.8,0);
 \draw[->,black!35] (0,-1.15) -- (0,1.2);
 \foreach \x/\y in {-0.8/-0.3,-0.55/0.42,0.12/-0.55,0.48/0.3,0.92/-0.15}
  \fill (\x,\y) circle (1.4pt);
 \draw[black!55] (0.48,0.3) -- (0.85,1.03);
 \node at (0.85,1.25)
  {$\FeatureCoordinates(\ComparisonPoint)$};
 \draw[black!55] (-0.8,-0.3) -- (-0.95,-0.9);
 \node at (-0.95,-1.12)
  {$\FeatureCoordinates(\IntegrationPoint)$};
\end{scope}
\begin{scope}[xshift=5.8cm]
 \begin{scope}[rotate=30]
  \filldraw[fill=black!4] (0,0) ellipse (1.45 and 0.88);
 \end{scope}
 \draw[->,black!35] (-1.8,0) -- (1.8,0);
 \draw[->,black!35] (0,-1.15) -- (0,1.2);
 \begin{scope}[rotate=30]
  \foreach \x/\y in {-0.8/-0.3,-0.55/0.42,0.12/-0.55,0.48/0.3,0.92/-0.15}
   \fill (\x,\y) circle (1.4pt);
  \coordinate (rotatedy) at (0.48,0.3);
  \coordinate (rotatedz) at (-0.8,-0.3);
 \end{scope}
 \draw[black!55] (rotatedy) -- (0.4,1.03);
 \node at (0.4,1.25)
  {$\OrthogonalChange\FeatureCoordinates(\ComparisonPoint)$};
 \draw[black!55] (rotatedz) -- (-1,-0.92);
 \node at (-1,-1.14)
  {$\OrthogonalChange\FeatureCoordinates(\IntegrationPoint)$};
\end{scope}
\draw[->] (2.05,0.15) -- node[above]
 {$\OrthogonalChange\in\OrthogonalGroup(2)$} (3.75,0.15);
\node at (0,-1.55)
 {$\{\FirstVector\mid\NormSymbol(\FirstVector)\leq1\}$};
\node at (5.8,-1.55)
 {$\{\FirstVector\mid
     (\NormSymbol\circ\OrthogonalChange^{-1})(\FirstVector)\leq1\}$};
\node at (2.9,-2.35)
 {$(\NormSymbol\circ\OrthogonalChange^{-1})
     \bigl(\OrthogonalChange\FeatureCoordinates(\ComparisonPoint)
           -\OrthogonalChange\FeatureCoordinates(\IntegrationPoint)\bigr)
   =\NormSymbol\bigl(\FeatureCoordinates(\ComparisonPoint)
                     -\FeatureCoordinates(\IntegrationPoint)\bigr)$};
\end{tikzpicture}
\caption{An orthogonal change of coordinates acts on both the coordinate
image and the norm unit ball.
The dots represent a finite coordinate image in dimension
$2$,
and the shaded regions represent the two unit balls.
Their simultaneous rotation leaves the induced pseudometric unchanged.}
\label{fig:orthogonal-coordinate-change}
\end{figure}


\begin{lemma}\label{lem:coordinate-pseudometric-estimate}
\paraid{p-29c8b5705912}
Let
$\BaseCarrier,\ComparisonCarrier$
be nonempty compact spaces,
let
$\Correspondence\subset\BaseCarrier\times\ComparisonCarrier$
be a correspondence,
and let
$\ModelDimension\geq1$.
Let
$\FeatureCoordinates_0\colon\BaseCarrier\to\RealNumbers^\ModelDimension$
and
$\FeatureCoordinates_1\colon\ComparisonCarrier\to\RealNumbers^\ModelDimension$
be continuous,
and let
$\NormSymbol_0,\NormSymbol_1$
be norms on
$\RealNumbers^\ModelDimension$.
Denote the induced pseudometrics by
$\Pseudometric_0$
and
$\Pseudometric_1$,
respectively.
Put
\[
 \CoefficientBound=\sup_{\BasePoint\in\BaseCarrier}
       \norm{\FeatureCoordinates_0(\BasePoint)}_2,
\]
\[
 \CoordinateComparisonError
 =\sup_{(\BasePoint,\ComparisonPoint)\in\Correspondence}
       \norm{\FeatureCoordinates_0(\BasePoint)-\FeatureCoordinates_1(\ComparisonPoint)}_2,
\]
\[
 \NormComparisonError
 =\sup_{\substack{\CoefficientVector\in\RealNumbers^\ModelDimension\\\norm{\CoefficientVector}_2=1}}
       |\NormSymbol_0(\CoefficientVector)-\NormSymbol_1(\CoefficientVector)|,
\]
and
\[
 \NormComparisonBound
 =\max_{\substack{\CoefficientVector\in\RealNumbers^\ModelDimension\\\norm{\CoefficientVector}_2=1}}\NormSymbol_1(\CoefficientVector).
\]
Then
\begin{equation}\label{eq:coordinate-pseudometric-estimate}
 \PseudometricError_\Correspondence(\Pseudometric_0,\Pseudometric_1)
 \leq2\CoefficientBound\NormComparisonError
       +2\NormComparisonBound\CoordinateComparisonError.
\end{equation}
\end{lemma}

\begin{proof}
\paraid{p-9d7ad7bd9f69}
Fix
$(\BasePoint_0,\ComparisonPoint_0),(\BasePoint_1,\ComparisonPoint_1)\in\Correspondence$
and set
\[
 \FirstVector=\FeatureCoordinates_0(\BasePoint_0)-\FeatureCoordinates_0(\BasePoint_1),
\]
and
\[
 \SecondVector=\FeatureCoordinates_1(\ComparisonPoint_0)-\FeatureCoordinates_1(\ComparisonPoint_1).
\]
Then
$\norm{\FirstVector}_2\leq2\CoefficientBound$
and
$\norm{\FirstVector-\SecondVector}_2\leq2\CoordinateComparisonError$.
Homogeneity and the norm triangle inequality imply
\[
 \begin{aligned}
 |\NormSymbol_0(\FirstVector)-\NormSymbol_1(\SecondVector)|
 &\leq|\NormSymbol_0(\FirstVector)-\NormSymbol_1(\FirstVector)|
       +\NormSymbol_1(\FirstVector-\SecondVector)\\
 &\leq\NormComparisonError\norm{\FirstVector}_2
       +\NormComparisonBound\norm{\FirstVector-\SecondVector}_2\\
 &\leq2\CoefficientBound\NormComparisonError
       +2\NormComparisonBound\CoordinateComparisonError.
 \end{aligned}
\]
Taking the supremum over the two pairs proves
\eqref{eq:coordinate-pseudometric-estimate}.
\end{proof}

\paraid{p-03fcd50a02a7}
When a function is defined on GH isometry classes, its value at
$\MetricSpace{\ComparisonCarrier}{\ComparisonMetric}$ denotes its value at
$[\ComparisonCarrier,\ComparisonMetric]$, the isometry class of that representative.
In particular, this convention applies to the error functions.

\begin{theorem}\label{thm:local-models}
\paraid{p-9559de6c1eef}
Let
$\MetricSpace{\BaseCarrier }{\MetricSymbol }\in\GHSpace$
be a nonempty compact metric space and let
$\LocalTolerance >0$.
Then there exist an open subset
$\ModelNeighborhood\subset\GHSpace$
with
$\MetricSpace{\BaseCarrier }{\MetricSymbol }\in\ModelNeighborhood$
and an integer
$\ModelDimension \geq1$
and a continuous function
$\LocalError\colon\ModelNeighborhood\to[0,\infty)$,
together with the following simultaneous choices.
For each concrete compact metric space
$\MetricSpace{\ComparisonCarrier}{\ComparisonMetric}$
whose isometry class belongs to
$\ModelNeighborhood$,
choose a class
\[
 \LocalModelClass{\ComparisonCarrier}{\ComparisonMetric}
 \in\CoordinateClassSpace{\ModelDimension}{\ComparisonCarrier}.
\]
These choices satisfy the following properties.
\begin{enumerate}[label=\textup{(A\arabic*)},ref=\textup{(A\arabic*)}]
\item\label{item:model-class}
\paraid{p-578e284641e7}
For each such compact metric space
$\MetricSpace{\ComparisonCarrier}{\ComparisonMetric}$,
the assigned class satisfies
\[
 \LocalModelClass{\ComparisonCarrier}{\ComparisonMetric}
 \in\CoordinateClassSpace{\ModelDimension}{\ComparisonCarrier}.
\]
In \ref{item:model-bound} and \ref{item:model-pseudometric}--\ref{item:model-center-error},
$(\FeatureCoordinates_\ComparisonCarrier,\NormSymbol_\ComparisonCarrier)$
denotes an arbitrary coordinate pair in this class.
The bound in \ref{item:model-bound} and the pseudometric and error assertions
in \ref{item:model-pseudometric}--\ref{item:model-center-error} hold for every coordinate pair in this class.
\item\label{item:model-equivariance}
\paraid{p-38960d79969a}
For every isometry
$\InputIsometry\colon\MetricSpace{\ComparisonCarrier}{\ComparisonMetric}
 \to\MetricSpace{\ComparisonCarrier_0}{\ComparisonMetric_0}$
between compact metric spaces whose isometry classes belong to
$\ModelNeighborhood$
and every choice of coordinate pairs
$(\FeatureCoordinates_\ComparisonCarrier,\NormSymbol_\ComparisonCarrier)$
and
$(\FeatureCoordinates_{\ComparisonCarrier_0},\NormSymbol_{\ComparisonCarrier_0})$
of the two assigned coordinate classes,
there exists
$\OrthogonalChange\in\OrthogonalGroup(\ModelDimension)$
such that
\begin{equation}\label{eq:input-isometry-1}
 \FeatureCoordinates_{\ComparisonCarrier_0}\circ\InputIsometry
 =\OrthogonalChange\circ\FeatureCoordinates_\ComparisonCarrier,
\end{equation}
and
\begin{equation}\label{eq:input-isometry-2}
 \NormSymbol_{\ComparisonCarrier_0}=\NormSymbol_\ComparisonCarrier\circ\OrthogonalChange^{-1}.
\end{equation}
\item\label{item:model-bound}
\paraid{p-ac13929fbe66}
For every
$\ComparisonPoint \in \ComparisonCarrier $,
\[
 \NormSymbol _\ComparisonCarrier (\FeatureCoordinates _\ComparisonCarrier (\ComparisonPoint ))\leq2\diam\MetricSpace{\ComparisonCarrier }{\ComparisonMetric }.
\]
\item\label{item:model-continuity}
\paraid{p-e36eef01289e}
Let
$\{\MetricSpace{\ComparisonCarrier_\SequenceIndex}{\ComparisonMetric_\SequenceIndex}\}_{\SequenceIndex\in\NonnegativeIntegers}$
be a sequence of compact metric spaces whose isometry classes belong to
$\ModelNeighborhood$,
and let
$\MetricSpace{\ComparisonCarrier}{\ComparisonMetric}$
be a compact metric space whose isometry class belongs to
$\ModelNeighborhood$.
Assume that these spaces have isometric embeddings into a compact
$\MetricSpace{\AmbientCarrier }{\AmbientMetric }$
with
$\HausdorffDistance{\AmbientMetric }(\ComparisonCarrier _\SequenceIndex ,\ComparisonCarrier )\to0$.
Then there exist a coordinate pair
$(\FeatureCoordinates_\ComparisonCarrier,\NormSymbol_\ComparisonCarrier)$
and a family of coordinate pairs
$\{(\FeatureCoordinates_{\ComparisonCarrier_\SequenceIndex},\NormSymbol_{\ComparisonCarrier_\SequenceIndex})\}_{\SequenceIndex\in\NonnegativeIntegers}$
such that
\[
 (\FeatureCoordinates_\ComparisonCarrier,\NormSymbol_\ComparisonCarrier)
 \in\LocalModelClass{\ComparisonCarrier}{\ComparisonMetric},
\]
and
\[
 (\FeatureCoordinates_{\ComparisonCarrier_\SequenceIndex},\NormSymbol_{\ComparisonCarrier_\SequenceIndex})
 \in\LocalModelClass{\ComparisonCarrier_\SequenceIndex}{\ComparisonMetric_\SequenceIndex},
\]
and
\[
 \sup_{\substack{\CoefficientVector\in\RealNumbers^\ModelDimension\\\norm{\CoefficientVector}_2=1}}|\NormSymbol _{\ComparisonCarrier _\SequenceIndex }(\CoefficientVector )-\NormSymbol _\ComparisonCarrier (\CoefficientVector )|\to0.
\]
These coordinate pairs can be chosen using only the prescribed sequence
and ambient embeddings, independently of any sequence of correspondences with vanishing ambient displacement.
For every sequence of correspondences
$\Correspondence _\SequenceIndex \subset \ComparisonCarrier _\SequenceIndex \times \ComparisonCarrier $
satisfying
\[
 \sup_{(\ComparisonPoint _\SequenceIndex ,\ComparisonPoint )\in \Correspondence _\SequenceIndex }\AmbientMetric (\ComparisonPoint _\SequenceIndex ,\ComparisonPoint )\to0,
\]
the coordinate maps
$\FeatureCoordinates_{\ComparisonCarrier_\SequenceIndex}$
and
$\FeatureCoordinates_\ComparisonCarrier$
satisfy
\begin{equation}\label{eq:local-coordinate-convergence}
 \sup_{(\ComparisonPoint_\SequenceIndex,\ComparisonPoint)\in\Correspondence_\SequenceIndex}
 \norm{\FeatureCoordinates_{\ComparisonCarrier_\SequenceIndex}(\ComparisonPoint_\SequenceIndex)
 -\FeatureCoordinates_\ComparisonCarrier(\ComparisonPoint)}_2\to0.
\end{equation}
\item\label{item:model-pseudometric}
\paraid{p-29f1150553e1}
For
$\ComparisonPoint,\IntegrationPoint\in\ComparisonCarrier$,
the formula
\begin{equation}\label{eq:local-model-pseudometric}
 \Pseudometric _\ComparisonCarrier (\ComparisonPoint ,\IntegrationPoint )=\NormSymbol _\ComparisonCarrier (\FeatureCoordinates _\ComparisonCarrier (\ComparisonPoint )-\FeatureCoordinates _\ComparisonCarrier (\IntegrationPoint ))
\end{equation}
defines a continuous pseudometric on
$\ComparisonCarrier $.
\item\label{item:model-pseudometric-invariance}
\paraid{p-b7eec021787e}
For every isometry
$\InputIsometry\colon\MetricSpace{\ComparisonCarrier}{\ComparisonMetric}
 \to\MetricSpace{\ComparisonCarrier_0}{\ComparisonMetric_0}$
between compact metric spaces whose isometry classes belong to
$\ModelNeighborhood$
and every
$\ComparisonPoint,\IntegrationPoint\in\ComparisonCarrier$,
\[
 \Pseudometric_{\ComparisonCarrier_0}(\InputIsometry(\ComparisonPoint),\InputIsometry(\IntegrationPoint))
 =\Pseudometric_\ComparisonCarrier(\ComparisonPoint,\IntegrationPoint).
\]
\item\label{item:model-error-definition}
\paraid{p-702b7e8de607}
For each compact metric representative
$\MetricSpace{\ComparisonCarrier}{\ComparisonMetric}$
of a class in
$\ModelNeighborhood$,
the error function satisfies
\[
 \LocalError \MetricSpace{\ComparisonCarrier }{\ComparisonMetric }=\sup_{\ComparisonPoint ,\IntegrationPoint \in \ComparisonCarrier }|\Pseudometric _\ComparisonCarrier (\ComparisonPoint ,\IntegrationPoint )-\ComparisonMetric (\ComparisonPoint ,\IntegrationPoint )|.
\]
By \ref{item:model-pseudometric-invariance}, this value is independent
of the representative and of the coordinate pair.
\item\label{item:model-error-continuity}
\paraid{p-685c2b6e8303}
The error function
$\LocalError$
is continuous.
\item\label{item:model-center-error}
\paraid{p-f46126c20706}
At the center, it satisfies
$\LocalError \MetricSpace{\BaseCarrier }{\MetricSymbol }<\LocalTolerance $.
\end{enumerate}
\end{theorem}

\paraid{p-fc5c3b7ac001}
Write
$\mathcal L$
for the choice rule that associates to each concrete compact metric space
$\MetricSpace{\ComparisonCarrier}{\ComparisonMetric}$
whose class belongs to
$\ModelNeighborhood$
the class
\[
 \mathcal L_{\ComparisonCarrier,\ComparisonMetric}
 =\LocalModelClass{\ComparisonCarrier}{\ComparisonMetric}
 \in\CoordinateClassSpace{\ModelDimension}{\ComparisonCarrier}.
\]
We call the data
\[
 \mathfrak M=(\ModelNeighborhood,\ModelDimension,\mathcal L,\LocalError)
\]
a \emph{local model}.
Condition \ref{item:model-equivariance} describes how these values are transported by isometries.
Condition \ref{item:model-continuity} asserts the existence of a convergent family of coordinate pairs
for each prescribed sequence and its ambient embeddings.
These pairs represent the classes of
$\mathfrak M$.
By conditions \ref{item:model-pseudometric-invariance}--\ref{item:model-error-continuity},
the rule
\[
 \MetricSpace{\ComparisonCarrier}{\ComparisonMetric}
 \longmapsto\MetricQuotient{\ComparisonCarrier}{\Pseudometric_\ComparisonCarrier}
\]
induces a well-defined map
$\ModelNeighborhood\to\GHSpace$.
The error function
$\LocalError$
is independently well defined on isometry classes by
\ref{item:model-error-definition}.
The map
$\FeatureCoordinates _\ComparisonCarrier $
need not be injective,
which is why
$\Pseudometric _\ComparisonCarrier $
is a pseudometric.

\begin{proof}
\paraid{p-6d8a3f3aa620}
For the singleton center and each
$\MetricSpace{\ComparisonCarrier}{\ComparisonMetric}\in\GHSpace$,
set
$\ModelDimension =1$,
$\NormSymbol _\ComparisonCarrier =|\cdot|$,
and
$\FeatureCoordinates _\ComparisonCarrier =0$.
Define
$\LocalModelClass{\ComparisonCarrier}{\ComparisonMetric}
 =[(0,|\cdot|)]_{\OrthogonalGroup(1)}$.
Then
$\Pseudometric _\ComparisonCarrier =0$
and
$\LocalError \MetricSpace{\ComparisonCarrier }{\ComparisonMetric }=\diam\MetricSpace{\ComparisonCarrier }{\ComparisonMetric }$.
The error is continuous because the diameter is continuous for
$\GHDistance$.
For a prescribed
$\LocalTolerance>0$,
we may take its domain to be the open neighborhood
\[
 \ModelNeighborhood
 =\{\MetricSpace{\ComparisonCarrier}{\ComparisonMetric}\in\GHSpace
       \mid\diam\MetricSpace{\ComparisonCarrier}{\ComparisonMetric}<\LocalTolerance\}
 \ni\SingletonSpace.
\]
On this domain the zero maps and fixed norm satisfy
\ref{item:model-bound}--\ref{item:model-center-error},
with the error equal to the diameter.
The coordinate maps
$\FeatureCoordinates_\ComparisonCarrier$
vanish and the norm
$\NormSymbol_\ComparisonCarrier$
is fixed,
so the coordinate and isometry identities also hold.
This proves the theorem when the center is a singleton.

\paraid{p-08d995d825ac}
Assume henceforth that the center is not a singleton.
We construct the local data from a single positive real number
$\SpectralCutoff$,
prove their continuity first on each carrier and then as the carrier varies,
and finally verify invariance and the error estimate.

\ProofStep{A neighborhood with constant spectral dimension.}\label{step:local-model-spectral-dimension}
\paraid{p-3a9b99318187}
Choose
$\Radius>0$
sufficiently small.
Then choose
$2\leq\IntegrabilityExponent<\infty$
sufficiently large.
Finally,
choose
$\SmallError>0$
sufficiently small and a corresponding positive real number
$\SpectralCutoff $
for
$\MetricSpace{\BaseCarrier }{\MetricSymbol }$
so that
\[
\begin{gathered}
 \norm{\Pseudometric_\BaseCarrier-\MetricSymbol}_\infty
 \leq2\SmallError
      +\DiameterBound(1-\BallMassBound_\Radius^{1/\IntegrabilityExponent})
      +2\Radius
 <\LocalTolerance,
\\
 \SpectralCutoff,-\SpectralCutoff
 \notin\spec(\DistanceOperator{\BaseCarrier}{\MetricSymbol}),
 \qquad
 \dim\SpectralSubspace{\BaseCarrier}{\SpectralCutoff}
 =\ModelDimension\geq1.
\end{gathered}
\]
Define
\[
 \ModelNeighborhood=
 \{\MetricSpace{\ComparisonCarrier }{\ComparisonMetric }\in\GHSpace\mid
     \SpectralCutoff ,-\SpectralCutoff \notin\spec(\DistanceOperator{\ComparisonCarrier}{\ComparisonMetric}),\
     \dim\SpectralSubspace{\ComparisonCarrier }{\SpectralCutoff }=\ModelDimension \}.
\]
Then
$\MetricSpace{\BaseCarrier}{\MetricSymbol}\in\ModelNeighborhood$.
We prove that
$\ModelNeighborhood$
is open.
Let
$\MetricSpace{\ComparisonCarrier_\SequenceIndex}{\ComparisonMetric_\SequenceIndex}$
converge to
$\MetricSpace{\ComparisonCarrier}{\ComparisonMetric}\in\ModelNeighborhood$
in
$\GHSpace$.
By \Cref{lem:common-gh-embedding},
we realize the sequence and its limit as compact subsets of a common compact metric space
$\MetricSpace{\AmbientCarrier}{\AmbientMetric}$
with
\[
 \HausdorffDistance{\AmbientMetric}
 (\ComparisonCarrier_\SequenceIndex,\ComparisonCarrier)\to0.
\]
\Cref{thm:invariant-measures} implies
\[
 \SelectedMeasure{\ComparisonCarrier_\SequenceIndex}{\ComparisonMetric_\SequenceIndex}
 \to\SelectedMeasure{\ComparisonCarrier}{\ComparisonMetric}
 \quad\text{weakly on }\AmbientCarrier.
\]
\Cref{lem:spectral-projections} now yields
$\MetricSpace{\ComparisonCarrier_\SequenceIndex}{\ComparisonMetric_\SequenceIndex}
 \in\ModelNeighborhood$
for all sufficiently large
$\SequenceIndex$.
Since
$\GHSpace$
is metrizable,
this proves openness.

\ProofStep{Norms and continuous coordinate maps.}\label{step:local-model-coordinates}
\paraid{p-5e1efa60ac95}
For
$\MetricSpace{\ComparisonCarrier }{\ComparisonMetric }\in\ModelNeighborhood$,
choose a real orthonormal basis
$\BasisFunction_0,\ldots,\BasisFunction_{\ModelDimension-1}$
of
$\SpectralSubspace{\ComparisonCarrier }{\SpectralCutoff }$.
Define the functional
$\NormSymbol_\ComparisonCarrier\colon\RealNumbers^\ModelDimension\to[0,\infty)$
as follows.
Use
$\IntegrationPoint\in\ComparisonCarrier$
as the integration variable in
\eqref{eq:local-norm} and \eqref{eq:best-approximation-objective}.
For every
$\CoefficientVector\in\RealNumbers^\ModelDimension$,
put
\begin{equation}\label{eq:local-norm}
 \NormSymbol _\ComparisonCarrier (\CoefficientVector )=\left(\int_\ComparisonCarrier \left|\sum_{\BasisIndex=0}^{\ModelDimension-1}  \CoefficientVector _\BasisIndex \BasisFunction _\BasisIndex (\IntegrationPoint )\right|^\IntegrabilityExponent
                                     \IntegrationDifferential \SelectedMeasure{\ComparisonCarrier}{\ComparisonMetric}(\IntegrationPoint )\right)^{1/\IntegrabilityExponent }.
\end{equation}
Since
$\IntegrabilityExponent \geq2$
and the measure
$\SelectedMeasure{\ComparisonCarrier}{\ComparisonMetric}$
is a probability,
orthonormality implies for every
$\CoefficientVector\in\RealNumbers^\ModelDimension$
\begin{equation}\label{eq:local-norm-euclidean-lower-bound}
 \NormSymbol _\ComparisonCarrier (\CoefficientVector )\geq\left\|\sum_{\BasisIndex=0}^{\ModelDimension-1}  \CoefficientVector _\BasisIndex \BasisFunction _\BasisIndex \right\|_2=\norm{\CoefficientVector }_2.
\end{equation}
The functional is homogeneous and satisfies the triangle inequality by the
$\LebesgueSpace^\IntegrabilityExponent$-norm properties.
The lower bound in \eqref{eq:local-norm-euclidean-lower-bound} proves nondegeneracy,
so
$\NormSymbol_\ComparisonCarrier$ is a norm.

\paraid{p-8cbc0bb85af6}
Define the objective function
$\ApproximationObjective_\ComparisonCarrier\colon\ComparisonCarrier\times\RealNumbers^\ModelDimension\to[0,\infty)$
as follows.
For every
$\ComparisonPoint\in\ComparisonCarrier$
and
$\CoefficientVector\in\RealNumbers^\ModelDimension$,
put
\begin{equation}\label{eq:best-approximation-objective}
 \ApproximationObjective _\ComparisonCarrier (\ComparisonPoint ,\CoefficientVector )
 =\int_\ComparisonCarrier \left|\ComparisonMetric (\ComparisonPoint ,\IntegrationPoint )-\sum_{\BasisIndex=0}^{\ModelDimension-1}  \CoefficientVector _\BasisIndex \BasisFunction _\BasisIndex (\IntegrationPoint )\right|^\IntegrabilityExponent
                                      \IntegrationDifferential \SelectedMeasure{\ComparisonCarrier}{\ComparisonMetric}(\IntegrationPoint ).
\end{equation}
For each
$\ComparisonPoint\in\ComparisonCarrier$,
define
$\FeatureCoordinates_\ComparisonCarrier(\ComparisonPoint)\in\RealNumbers^\ModelDimension$
to be the unique minimizer of
$\CoefficientVector\mapsto\ApproximationObjective_\ComparisonCarrier(\ComparisonPoint,\CoefficientVector)$.
These minimizers define a map
$\FeatureCoordinates_\ComparisonCarrier\colon\ComparisonCarrier\to\RealNumbers^\ModelDimension$.
Existence and uniqueness follow from
\Cref{thm:best-approximation},
applied to the finite-dimensional spectral subspace
$\SpectralSubspace{\ComparisonCarrier}{\SpectralCutoff}$,
viewed as a subspace of
$\LebesgueSpace^\IntegrabilityExponent(\ComparisonCarrier,
\SelectedMeasure{\ComparisonCarrier}{\ComparisonMetric})$.
For each
$\ComparisonPoint\in\ComparisonCarrier$,
define the function
$\BestApproximationFunction_{\ComparisonCarrier,\ComparisonPoint}\colon\ComparisonCarrier\to\RealNumbers$
by setting, for every
$\IntegrationPoint\in\ComparisonCarrier$,
\begin{equation}\label{eq:local-best-approximation-function}
 \BestApproximationFunction_{\ComparisonCarrier,\ComparisonPoint}(\IntegrationPoint)
 =\sum_{\BasisIndex=0}^{\ModelDimension-1}
 (\FeatureCoordinates_\ComparisonCarrier(\ComparisonPoint))_\BasisIndex
 \BasisFunction_\BasisIndex(\IntegrationPoint).
\end{equation}
By \eqref{eq:best-approximation-objective},
$\BestApproximationFunction_{\ComparisonCarrier,\ComparisonPoint}$
is the best approximation of
$\ComparisonMetric(\ComparisonPoint,\cdot)$
in
$\SpectralSubspace{\ComparisonCarrier}{\SpectralCutoff}$.
The components of the vector
$\FeatureCoordinates_\ComparisonCarrier(\ComparisonPoint)$
are the coefficients of this best approximation in the chosen orthonormal basis
$\{\BasisFunction_\BasisIndex\}_{0\leq\BasisIndex<\ModelDimension}$.
Figure~\ref{fig:local-coordinates} displays this construction of
$\FeatureCoordinates_\ComparisonCarrier$.
\begin{figure}[H]
\centering
\begin{tikzpicture}[>=Stealth,node distance=12mm,font=\small]
\node (profile) {$\ComparisonMetric(\ComparisonPoint,\cdot)
 \in\LebesgueSpace^\IntegrabilityExponent
 (\ComparisonCarrier,\SelectedMeasure{\ComparisonCarrier}{\ComparisonMetric})$};
\node[below=of profile] (approximation) {$
 \displaystyle\sum_{\BasisIndex=0}^{\ModelDimension-1}
 (\FeatureCoordinates_\ComparisonCarrier(\ComparisonPoint))_\BasisIndex
 \BasisFunction_\BasisIndex
 \in\SpectralSubspace{\ComparisonCarrier}{\SpectralCutoff}$};
\node[below=of approximation] (coordinates) {$
 \FeatureCoordinates_\ComparisonCarrier(\ComparisonPoint)
 \in(\RealNumbers^\ModelDimension,\NormSymbol_\ComparisonCarrier)$};
\draw[->] (profile) -- node[right]
 {best approximation in $\LebesgueSpace^\IntegrabilityExponent$} (approximation);
\draw[->] (approximation) -- node[right]
 {coordinates in $\{\BasisFunction_\BasisIndex\}_{0\leq\BasisIndex<\ModelDimension}$} (coordinates);
\end{tikzpicture}
\caption{The best approximation of
$\ComparisonMetric(\ComparisonPoint,\cdot)$
in
$\SpectralSubspace{\ComparisonCarrier}{\SpectralCutoff}$
has coordinate vector
$\FeatureCoordinates_\ComparisonCarrier(\ComparisonPoint)$
in the chosen
$\LebesgueSpace^2$-orthonormal
basis.
This vector is the minimizer in \eqref{eq:best-approximation-objective}.
With the
$\LebesgueSpace^\IntegrabilityExponent$
norm on
$\SpectralSubspace{\ComparisonCarrier}{\SpectralCutoff}$,
the coordinate map in the second arrow is a linear isometry onto
$(\RealNumbers^\ModelDimension,\NormSymbol_\ComparisonCarrier)$
by \eqref{eq:local-norm}.}
\label{fig:local-coordinates}
\end{figure}


\paraid{p-aa449f741605}
For each
$\ComparisonPoint\in\ComparisonCarrier$,
write
$\ComparisonMetric_\ComparisonPoint=\ComparisonMetric(\ComparisonPoint,\cdot)
 \in\ContinuousFunctions(\ComparisonCarrier)$
for the distance profile at
$\ComparisonPoint$.
Since
$\SelectedMeasure{\ComparisonCarrier}{\ComparisonMetric}$
is a probability and
$\ComparisonMetric(\ComparisonPoint,\IntegrationPoint)\leq
\diam\MetricSpace{\ComparisonCarrier}{\ComparisonMetric}$
for every
$\IntegrationPoint\in\ComparisonCarrier$,
we have
\[
 \norm{\ComparisonMetric_\ComparisonPoint}_\IntegrabilityExponent^{\IntegrabilityExponent}
 =\int_\ComparisonCarrier
   |\ComparisonMetric(\ComparisonPoint,\IntegrationPoint)|^{\IntegrabilityExponent}
   \,\IntegrationDifferential\SelectedMeasure{\ComparisonCarrier}{\ComparisonMetric}(\IntegrationPoint)
 \leq\diam\MetricSpace{\ComparisonCarrier}{\ComparisonMetric}^{\IntegrabilityExponent}
   \int_\ComparisonCarrier
   \,\IntegrationDifferential\SelectedMeasure{\ComparisonCarrier}{\ComparisonMetric}(\IntegrationPoint)
 =\diam\MetricSpace{\ComparisonCarrier}{\ComparisonMetric}^{\IntegrabilityExponent}.
\]
By minimality of
$\FeatureCoordinates_\ComparisonCarrier(\ComparisonPoint)$,
\[
 \norm{\ComparisonMetric_\ComparisonPoint
       -\BestApproximationFunction_{\ComparisonCarrier,\ComparisonPoint}}_\IntegrabilityExponent
 \leq\norm{\ComparisonMetric_\ComparisonPoint}_\IntegrabilityExponent.
\]
Therefore, the triangle inequality and the definition of
$\NormSymbol_\ComparisonCarrier$
yield
\[
 \NormSymbol_\ComparisonCarrier
   (\FeatureCoordinates_\ComparisonCarrier(\ComparisonPoint))
 =\norm{\BestApproximationFunction_{\ComparisonCarrier,\ComparisonPoint}}_\IntegrabilityExponent
 \leq\norm{\ComparisonMetric_\ComparisonPoint}_\IntegrabilityExponent
    +\norm{\ComparisonMetric_\ComparisonPoint
           -\BestApproximationFunction_{\ComparisonCarrier,\ComparisonPoint}}_\IntegrabilityExponent
 \leq2\diam\MetricSpace{\ComparisonCarrier}{\ComparisonMetric}.
\]
Together with the lower bound in
\eqref{eq:local-norm-euclidean-lower-bound},
this proves
\begin{equation}\label{eq:coefficient-bounds}
 \norm{\FeatureCoordinates _\ComparisonCarrier (\ComparisonPoint )}_2\leq \NormSymbol _\ComparisonCarrier (\FeatureCoordinates _\ComparisonCarrier (\ComparisonPoint ))\leq2\diam\MetricSpace{\ComparisonCarrier }{\ComparisonMetric }.
\end{equation}
This proves \ref{item:model-bound}.
In particular,
the minimizers stay in a fixed compact Euclidean ball when the diameters
are bounded.

\paraid{p-29af57cad7b7}
For
$\CoefficientBound>0$,
write
\[
 \CoefficientBall_\CoefficientBound
 =\{\CoefficientVector\in\RealNumbers^\ModelDimension
      \mid\norm{\CoefficientVector}_2\leq\CoefficientBound\}.
\]
For fixed
$\ComparisonCarrier$,
the function
$\ApproximationObjective_\ComparisonCarrier$
is continuous.
If
$\ComparisonPoint_\SubsequenceIndex\to\ComparisonPoint$,
continuity on the product of the compact space
$\ComparisonCarrier$
and each
$\CoefficientBall_\CoefficientBound$
implies
\[
 \sup_{\CoefficientVector\in\CoefficientBall_\CoefficientBound}
 |\ApproximationObjective_\ComparisonCarrier(\ComparisonPoint_\SubsequenceIndex,\CoefficientVector)
  -\ApproximationObjective_\ComparisonCarrier(\ComparisonPoint,\CoefficientVector)|\to0.
\]
The bound \eqref{eq:coefficient-bounds} and
\Cref{lem:minimizer-continuity} prove continuity of
$\FeatureCoordinates_\ComparisonCarrier$.

\ProofStep{Convergence of the objectives and norms.}\label{step:local-model-norm-convergence}
\paraid{p-4cd36224c006}
Assume that
$\HausdorffDistance{\AmbientMetric}(\ComparisonCarrier_\SequenceIndex,\ComparisonCarrier)\to0$
in the stated compact ambient space.
Write
$\ComparisonMetric_\SequenceIndex=\AmbientMetric|_{\ComparisonCarrier_\SequenceIndex\times\ComparisonCarrier_\SequenceIndex}$
and
$\ComparisonMetric=\AmbientMetric|_{\ComparisonCarrier\times\ComparisonCarrier}$.
Let $\iota_\SequenceIndex\colon\ComparisonCarrier_\SequenceIndex\hookrightarrow\AmbientCarrier$
and $\iota\colon\ComparisonCarrier\hookrightarrow\AmbientCarrier$ be the inclusions.
Define $\bar\mu_\SequenceIndex=(\iota_\SequenceIndex)_*\SelectedMeasure{\ComparisonCarrier_\SequenceIndex}{\ComparisonMetric_\SequenceIndex}$
and $\bar\mu=\iota_*\SelectedMeasure{\ComparisonCarrier}{\ComparisonMetric}$.
By \Cref{thm:invariant-measures},
\[
 \bar\mu_\SequenceIndex\to\bar\mu
 \quad\text{weakly on }\AmbientCarrier.
\]
Apply \Cref{lem:spectral-continuity} to the basis
$\{\BasisFunction_\BasisIndex\}_{0\leq\BasisIndex<\ModelDimension}$
of
$\SpectralSubspace{\ComparisonCarrier}{\SpectralCutoff}$
chosen in Step~\ref{step:local-model-coordinates}.
We obtain real orthonormal bases
$\{\BasisFunction_{\SequenceIndex,\BasisIndex}\}_{0\leq\BasisIndex<\ModelDimension}$
of
$\SpectralSubspace{\ComparisonCarrier_\SequenceIndex}{\SpectralCutoff}$
and
continuous functions
$\ExtendedBasisFunction_{\SequenceIndex,\BasisIndex},\ExtendedBasisFunction_\BasisIndex
 \colon\AmbientCarrier\to\RealNumbers$
such that
\[
 \ExtendedBasisFunction_{\SequenceIndex,\BasisIndex}|_{\ComparisonCarrier_\SequenceIndex}
 =\BasisFunction_{\SequenceIndex,\BasisIndex},
\]
and
\[
 \ExtendedBasisFunction_\BasisIndex|_\ComparisonCarrier=\BasisFunction_\BasisIndex,
\]
and
\[
 \max_{0\leq\BasisIndex<\ModelDimension}
 \norm{\ExtendedBasisFunction_{\SequenceIndex,\BasisIndex}-\ExtendedBasisFunction_\BasisIndex}_\infty\to0.
\]
Define
$\ApproximationObjective_\SequenceIndex,\ApproximationObjective
 \colon\AmbientCarrier\times\RealNumbers^\ModelDimension\to[0,\infty)$
as follows.
In both integrals below,
$\SecondVector$
ranges over
$\AmbientCarrier$.
For every sufficiently large
$\SequenceIndex\in\NonnegativeIntegers$,
every
$\IntegrationPoint\in\AmbientCarrier$,
and every
$\CoefficientVector\in\RealNumbers^\ModelDimension$,
put
\begin{equation}\label{eq:ambient-approximation-objectives-1}
 \ApproximationObjective _\SequenceIndex (\IntegrationPoint ,\CoefficientVector )
 =\int_\AmbientCarrier \left|\AmbientMetric (\IntegrationPoint ,\SecondVector )-\sum_{\BasisIndex=0}^{\ModelDimension-1}  \CoefficientVector _\BasisIndex \ExtendedBasisFunction _{\SequenceIndex ,\BasisIndex }(\SecondVector )\right|^\IntegrabilityExponent
                                     \IntegrationDifferential \bar\mu_\SequenceIndex(\SecondVector ).
\end{equation}
For every
$\IntegrationPoint\in\AmbientCarrier$
and
$\CoefficientVector\in\RealNumbers^\ModelDimension$,
put
\begin{equation}\label{eq:ambient-approximation-objectives-2}
 \ApproximationObjective(\IntegrationPoint,\CoefficientVector)
 =\int_\AmbientCarrier\left|\AmbientMetric(\IntegrationPoint,\SecondVector)
       -\sum_{\BasisIndex=0}^{\ModelDimension-1}
        \CoefficientVector_\BasisIndex\ExtendedBasisFunction_\BasisIndex(\SecondVector)
       \right|^\IntegrabilityExponent
       \IntegrationDifferential\bar\mu(\SecondVector).
\end{equation}
By \eqref{eq:best-approximation-objective},
\[
 \ApproximationObjective_\SequenceIndex
 |_{\ComparisonCarrier_\SequenceIndex\times\RealNumbers^\ModelDimension}
 =\ApproximationObjective_{\ComparisonCarrier_\SequenceIndex},
\]
and
\[
 \ApproximationObjective|_{\ComparisonCarrier\times\RealNumbers^\ModelDimension}
 =\ApproximationObjective_\ComparisonCarrier.
\]
For every
$\CoefficientBound >0$,
\begin{equation}\label{eq:uniform-objectives}
 \sup_{\substack{\IntegrationPoint\in\AmbientCarrier\\\CoefficientVector\in\RealNumbers^\ModelDimension\\\norm{\CoefficientVector}_2\leq\CoefficientBound}}
       |\ApproximationObjective _\SequenceIndex (\IntegrationPoint ,\CoefficientVector )-\ApproximationObjective (\IntegrationPoint ,\CoefficientVector )|\to0.
\end{equation}
Apply \Cref{lem:parameter-integrals} with compact parameter space
$\AmbientCarrier\times\{\CoefficientVector\in\RealNumbers^\ModelDimension
 \mid\norm{\CoefficientVector}_2\leq\CoefficientBound\}$.
Uniform convergence of the extended basis functions
$\ExtendedBasisFunction_{\SequenceIndex,\BasisIndex}$
to
$\ExtendedBasisFunction_\BasisIndex$
implies uniform convergence of the integrands in \eqref{eq:ambient-approximation-objectives-1} and \eqref{eq:ambient-approximation-objectives-2}
on this product with
$\AmbientCarrier$.
For the norms
$\NormSymbol_{\ComparisonCarrier_\SequenceIndex},\NormSymbol_\ComparisonCarrier$,
apply the same lemma with the Euclidean unit sphere
$\UnitSphere^{\ModelDimension-1}$
as parameter space and integrands
$\left|\sum_{\BasisIndex=0}^{\ModelDimension-1}\CoefficientVector_\BasisIndex
 \ExtendedBasisFunction_{\SequenceIndex,\BasisIndex}(\SecondVector)\right|^\IntegrabilityExponent$.
It follows that
\begin{equation}\label{eq:local-norm-power-convergence}
 \sup_{\substack{\CoefficientVector\in\RealNumbers^\ModelDimension\\\norm{\CoefficientVector}_2=1}}
 |\NormSymbol_{\ComparisonCarrier_\SequenceIndex}(\CoefficientVector)^\IntegrabilityExponent
 -\NormSymbol_\ComparisonCarrier(\CoefficientVector)^\IntegrabilityExponent|\to0.
\end{equation}
For nonnegative real numbers
$\FirstLength,\SecondLength$
and
$\IntegrabilityExponent\geq1$,
\begin{equation}\label{eq:root-holder-bound}
 |\FirstLength^{1/\IntegrabilityExponent}-\SecondLength^{1/\IntegrabilityExponent}|
 \leq|\FirstLength-\SecondLength|^{1/\IntegrabilityExponent}.
\end{equation}
Apply \eqref{eq:root-holder-bound} with
$\FirstLength=\NormSymbol_{\ComparisonCarrier_\SequenceIndex}(\CoefficientVector)^\IntegrabilityExponent$
and
$\SecondLength=\NormSymbol_\ComparisonCarrier(\CoefficientVector)^\IntegrabilityExponent$.
Taking the supremum and using \eqref{eq:local-norm-power-convergence},
we obtain
\begin{equation}\label{eq:local-norm-convergence}
 \begin{aligned}
 &\sup_{\substack{\CoefficientVector\in\RealNumbers^\ModelDimension\\\norm{\CoefficientVector}_2=1}}
 |\NormSymbol_{\ComparisonCarrier_\SequenceIndex}(\CoefficientVector)
       -\NormSymbol_\ComparisonCarrier(\CoefficientVector)|\\
 &\quad\leq
 \left(\sup_{\substack{\CoefficientVector\in\RealNumbers^\ModelDimension\\\norm{\CoefficientVector}_2=1}}
 |\NormSymbol_{\ComparisonCarrier_\SequenceIndex}(\CoefficientVector)^\IntegrabilityExponent
       -\NormSymbol_\ComparisonCarrier(\CoefficientVector)^\IntegrabilityExponent|
 \right)^{1/\IntegrabilityExponent}
 \longrightarrow0.
 \end{aligned}
\end{equation}

\ProofStep{Uniform convergence of the coordinate maps.}\label{step:local-model-coordinate-convergence}
\paraid{p-5f47a359eb40}
We now prove the second convergence in \ref{item:model-continuity}
for the bases
$\{\BasisFunction_{\SequenceIndex,\BasisIndex}\}_{0\leq\BasisIndex<\ModelDimension}$
and
$\{\BasisFunction_\BasisIndex\}_{0\leq\BasisIndex<\ModelDimension}$
chosen in Step~\ref{step:local-model-norm-convergence}.
Fix an arbitrary sequence of correspondences
$\{\Correspondence_\SequenceIndex\}_{\SequenceIndex\in\NonnegativeIntegers}$,
where
$\Correspondence_\SequenceIndex\subset\ComparisonCarrier_\SequenceIndex\times\ComparisonCarrier$
for every
$\SequenceIndex\in\NonnegativeIntegers$,
satisfying
\[
 \sup_{(\ComparisonPoint_\SequenceIndex,\ComparisonPoint)\in\Correspondence_\SequenceIndex}
 \AmbientMetric(\ComparisonPoint_\SequenceIndex,\ComparisonPoint)\to0.
\]
For the sake of contradiction,
suppose that \eqref{eq:local-coordinate-convergence} fails.
Choose
$\ErrorTolerance>0$,
a strictly increasing sequence of indices
$n_k$,
and for each
$k\in\NonnegativeIntegers$,
a pair
$(\ComparisonPoint _{n_k} ,\ComparisonPoint _{0,{n_k} })\in \Correspondence _{n_k} $
with
\begin{equation}\label{eq:nonconvergent-coordinate-separation}
 \norm{\FeatureCoordinates _{\ComparisonCarrier _{n_k} }(\ComparisonPoint _{n_k} )-\FeatureCoordinates _\ComparisonCarrier (\ComparisonPoint _{0,{n_k} })}_2\geq\ErrorTolerance .
\end{equation}
By compactness of
$\ComparisonCarrier$
and \eqref{eq:coefficient-bounds},
pass to a further subsequence and denote its indices again by
$n_k$.
For some
$\ComparisonPoint\in\ComparisonCarrier$
and
$\CoefficientVector\in\RealNumbers^\ModelDimension$,
we then obtain
\[
 \ComparisonPoint_{0,{n_k}}\to\ComparisonPoint,
\]
and
\[
 \FeatureCoordinates_{\ComparisonCarrier_{n_k}}(\ComparisonPoint_{n_k})
 \to\CoefficientVector.
\]
The condition on the correspondences implies
\[
 \AmbientMetric(\ComparisonPoint_{n_k},\ComparisonPoint)
 \leq\AmbientMetric(\ComparisonPoint_{n_k},\ComparisonPoint_{0,{n_k}})
      +\AmbientMetric(\ComparisonPoint_{0,{n_k}},\ComparisonPoint)
 \longrightarrow0.
\]
Fix an arbitrary
$\CoefficientBound>0$.
For every
$k\in\NonnegativeIntegers$,
define the uniform objective error by
\begin{equation}\label{eq:local-uniform-objective-error}
 a_k=\sup_{\substack{\IntegrationPoint\in\AmbientCarrier\\
                  \CoefficientVector\in\CoefficientBall_\CoefficientBound}}
 |\ApproximationObjective_{n_k}(\IntegrationPoint,\CoefficientVector)
  -\ApproximationObjective(\IntegrationPoint,\CoefficientVector)|.
\end{equation}
For every
$k\in\NonnegativeIntegers$,
define the error due to the moving evaluation point by
\begin{equation}\label{eq:local-moving-point-error}
 b_k=\sup_{\CoefficientVector\in\CoefficientBall_\CoefficientBound}
 |\ApproximationObjective(\ComparisonPoint_{n_k},\CoefficientVector)
  -\ApproximationObjective(\ComparisonPoint,\CoefficientVector)|.
\end{equation}
By \eqref{eq:uniform-objectives},
$a_k\to0$.
Uniform continuity of
$\ApproximationObjective$
on the compact set
$\AmbientCarrier\times\CoefficientBall_\CoefficientBound$
and convergence of
$\ComparisonPoint_{n_k}$
to
$\ComparisonPoint$
imply
$b_k\to0$.
Thus the triangle inequality yields
\begin{equation}\label{eq:local-moving-objective-convergence}
 \sup_{\CoefficientVector\in\CoefficientBall_\CoefficientBound}
 |\ApproximationObjective_{n_k}(\ComparisonPoint_{n_k},\CoefficientVector)
  -\ApproximationObjective(\ComparisonPoint,\CoefficientVector)|
 \leq a_k+b_k\longrightarrow0.
\end{equation}
The convergence in \eqref{eq:local-moving-objective-convergence} holds for every fixed
$\CoefficientBound>0$,
so the objectives converge uniformly on every compact subset of
$\RealNumbers^\ModelDimension$.
For every $k$, the vector
$\FeatureCoordinates_{\ComparisonCarrier_{n_k}}(\ComparisonPoint_{n_k})$
minimizes $\ApproximationObjective_{n_k}(\ComparisonPoint_{n_k},\cdot)$.
By \eqref{eq:coefficient-bounds},
choose
$\CoefficientBound>0$
large enough that
$\CoefficientBall_\CoefficientBound$
contains all these minimizers.
Equation \eqref{eq:local-moving-objective-convergence} proves uniform convergence of the objectives on that ball.
The limit objective has the unique minimizer
$\FeatureCoordinates_\ComparisonCarrier(\ComparisonPoint)$.
\Cref{lem:minimizer-continuity} therefore implies
$\CoefficientVector=\FeatureCoordinates_\ComparisonCarrier(\ComparisonPoint)$.
But continuity of
$\FeatureCoordinates _\ComparisonCarrier $
also implies
$\FeatureCoordinates _\ComparisonCarrier (\ComparisonPoint _{0,{n_k} })\to \FeatureCoordinates _\ComparisonCarrier (\ComparisonPoint )$,
contradicting \eqref{eq:nonconvergent-coordinate-separation}.
This proves \ref{item:model-continuity}.

\ProofStep{Changes of basis and isometries.}\label{step:local-model-isometries}
\paraid{p-67fdf3abc74d}
Let
$\OrthogonalChange\in\OrthogonalGroup(\ModelDimension)$
and let
$\{\widehat\BasisFunction_\BasisIndex\}_{0\leq\BasisIndex<\ModelDimension}$
be a second orthonormal basis such that for every
$0\leq\BasisIndex<\ModelDimension$
we have
\[
 \widehat\BasisFunction_\BasisIndex
 =\sum_{\SubsequenceIndex=0}^{\ModelDimension-1}
   \OrthogonalChange_{\BasisIndex\SubsequenceIndex}\BasisFunction_\SubsequenceIndex.
\]
For every
$\ComparisonPoint\in\ComparisonCarrier$
and
$\CoefficientVector\in\RealNumbers^\ModelDimension$,
the coordinate map
$\widehat\FeatureCoordinates_\ComparisonCarrier$
and the norm
$\widehat\NormSymbol_\ComparisonCarrier$
defined using the second basis satisfy
\[
 \widehat \FeatureCoordinates _\ComparisonCarrier (\ComparisonPoint )=\OrthogonalChange  \FeatureCoordinates _\ComparisonCarrier (\ComparisonPoint ),
\]
and
\[
 \widehat \NormSymbol _\ComparisonCarrier (\CoefficientVector )=\NormSymbol _\ComparisonCarrier (\OrthogonalChange ^{-1}\CoefficientVector ).
\]
Consequently,
$[\FeatureCoordinates_\ComparisonCarrier,\NormSymbol_\ComparisonCarrier]_{\OrthogonalGroup(\ModelDimension)}$
is independent of the orthonormal basis used in Step~\ref{step:local-model-coordinates}.
Define
\[
 \LocalModelClass{\ComparisonCarrier}{\ComparisonMetric}
 =[\FeatureCoordinates_\ComparisonCarrier,\NormSymbol_\ComparisonCarrier]_{\OrthogonalGroup(\ModelDimension)}.
\]
The coordinate pairs
$(\FeatureCoordinates_{\ComparisonCarrier_\SequenceIndex},
  \NormSymbol_{\ComparisonCarrier_\SequenceIndex})$
and
$(\FeatureCoordinates_\ComparisonCarrier,\NormSymbol_\ComparisonCarrier)$
constructed using the bases from Step~\ref{step:local-model-norm-convergence} represent
$\LocalModelClass{\ComparisonCarrier_\SequenceIndex}{\ComparisonMetric_\SequenceIndex}$
and
$\LocalModelClass{\ComparisonCarrier}{\ComparisonMetric}$,
respectively.
Their convergence in \eqref{eq:local-norm-convergence}
and \eqref{eq:local-coordinate-convergence}
proves \ref{item:model-continuity}.
For an isometry
$\InputIsometry \colon\MetricSpace{\ComparisonCarrier }{\ComparisonMetric }\to\MetricSpace{\ComparisonCarrier _0}{\ComparisonMetric _0}$,
define the pullback map
\[
 \IsometryPullback\colon
 \LebesgueSpace^2(\ComparisonCarrier_0,\SelectedMeasure{\ComparisonCarrier_0}{\ComparisonMetric_0})
 \to
 \LebesgueSpace^2(\ComparisonCarrier,\SelectedMeasure{\ComparisonCarrier}{\ComparisonMetric})
\]
as follows.
For every
$\HilbertFunction\in\LebesgueSpace^2(\ComparisonCarrier_0,\SelectedMeasure{\ComparisonCarrier_0}{\ComparisonMetric_0})$,
set
\[
 \IsometryPullback\HilbertFunction=\HilbertFunction\circ\InputIsometry.
\]
The map
$\IsometryPullback$
is unitary because
$\InputIsometry \Pushforward\SelectedMeasure{\ComparisonCarrier}{\ComparisonMetric}=\SelectedMeasure{\ComparisonCarrier _0}{\ComparisonMetric _0}$.
It intertwines the distance operators,
\[
 \DistanceOperator{\ComparisonCarrier}{\ComparisonMetric}(\HilbertFunction \circ \InputIsometry )
       =(\DistanceOperator{\ComparisonCarrier _0}{\ComparisonMetric _0}\HilbertFunction )\circ \InputIsometry .
\]
Thus
\[
 \IsometryPullback\bigl(\SpectralSubspace{\ComparisonCarrier_0}{\SpectralCutoff}\bigr)
 =\SpectralSubspace{\ComparisonCarrier}{\SpectralCutoff}.
\]
Let
$\{\BasisFunction_{0,\BasisIndex}\}_{0\leq\BasisIndex<\ModelDimension}$
be the orthonormal basis chosen on
$\ComparisonCarrier_0$.
Since both
$\{\BasisFunction_{0,\BasisIndex}\circ\InputIsometry\}_{\BasisIndex}$
and
$\{\BasisFunction_\BasisIndex\}_{\BasisIndex}$
are orthonormal bases of
$\SpectralSubspace{\ComparisonCarrier}{\SpectralCutoff}$,
there is
$\OrthogonalChange\in\OrthogonalGroup(\ModelDimension)$
such that for every
$0\leq\BasisIndex<\ModelDimension$
we have
\begin{equation}\label{eq:isometry-pullback-basis}
 \BasisFunction_{0,\BasisIndex}\circ\InputIsometry
 =\sum_{\SecondBasisIndex=0}^{\ModelDimension-1}
   \OrthogonalChange_{\BasisIndex\SecondBasisIndex}\BasisFunction_\SecondBasisIndex.
\end{equation}
Using \eqref{eq:isometry-pullback-basis} and
$\InputIsometry\Pushforward\SelectedMeasure{\ComparisonCarrier}{\ComparisonMetric}
 =\SelectedMeasure{\ComparisonCarrier_0}{\ComparisonMetric_0}$,
we obtain,
for every
$\ComparisonPoint\in\ComparisonCarrier$
and
$\CoefficientVector\in\RealNumbers^\ModelDimension$,
\begin{equation}\label{eq:isometry-objective-transform-1}
 \NormSymbol_{\ComparisonCarrier_0}(\CoefficientVector)
 =\NormSymbol_\ComparisonCarrier(\OrthogonalChange^{-1}\CoefficientVector),
\end{equation}
and
\begin{equation}\label{eq:isometry-objective-transform-2}
 \ApproximationObjective_{\ComparisonCarrier_0}
   (\InputIsometry(\ComparisonPoint),\CoefficientVector)
 =\ApproximationObjective_\ComparisonCarrier
   (\ComparisonPoint,\OrthogonalChange^{-1}\CoefficientVector).
\end{equation}
The identity \eqref{eq:isometry-objective-transform-2}
and uniqueness of the minimizers imply
\[
 \FeatureCoordinates_{\ComparisonCarrier_0}(\InputIsometry(\ComparisonPoint))
 =\OrthogonalChange\FeatureCoordinates_\ComparisonCarrier(\ComparisonPoint).
\]
This proves both identities in \eqref{eq:input-isometry-1} and \eqref{eq:input-isometry-2}.
Consequently,
for every
$\ComparisonPoint,\IntegrationPoint\in\ComparisonCarrier$,
\[
 \begin{aligned}
 \Pseudometric_{\ComparisonCarrier_0}
   (\InputIsometry(\ComparisonPoint),\InputIsometry(\IntegrationPoint))
 &=\NormSymbol_{\ComparisonCarrier_0}
   \bigl(\OrthogonalChange(\FeatureCoordinates_\ComparisonCarrier(\ComparisonPoint)
                         -\FeatureCoordinates_\ComparisonCarrier(\IntegrationPoint))\bigr)\\
 &=\Pseudometric_\ComparisonCarrier(\ComparisonPoint,\IntegrationPoint).
 \end{aligned}
\]
Thus the pseudometric
$\Pseudometric_\ComparisonCarrier$
is independent of the coordinates and is invariant under isometries of
$\MetricSpace{\ComparisonCarrier}{\ComparisonMetric}$.

\ProofStep{Continuity of the approximation error.}\label{step:local-model-error-continuity}
\paraid{p-2d2c5d189dec}
Equation \eqref{eq:local-model-pseudometric} defines a continuous pseudometric
because it pulls back the distance of a normed space by the continuous map
$\FeatureCoordinates_\ComparisonCarrier$.
Step~\ref{step:local-model-isometries} proves its invariance.
Set
\[
 \NormComparisonError_\SequenceIndex
 =\sup_{\substack{\CoefficientVector\in\RealNumbers^\ModelDimension\\\norm{\CoefficientVector}_2=1}}
       |\NormSymbol_{\ComparisonCarrier_\SequenceIndex}(\CoefficientVector)
          -\NormSymbol_\ComparisonCarrier(\CoefficientVector)|,
\]
and
\[
 \CoordinateComparisonError_\SequenceIndex
 =\sup_{(\ComparisonPoint_\SequenceIndex,\ComparisonPoint)\in\Correspondence_\SequenceIndex}
       \norm{\FeatureCoordinates_{\ComparisonCarrier_\SequenceIndex}(\ComparisonPoint_\SequenceIndex)
               -\FeatureCoordinates_\ComparisonCarrier(\ComparisonPoint)}_2.
\]
Steps~\ref{step:local-model-norm-convergence} and~\ref{step:local-model-coordinate-convergence} show that both numbers tend to
$0$.
Choose
$\CoefficientBound>0$
which bounds all the coordinate norms in \eqref{eq:coefficient-bounds}
along the sequence and on the limit space.
Such a bound exists because the diameters converge.
By \Cref{lem:coordinate-pseudometric-estimate},
\begin{equation}\label{eq:local-norm-pair-estimate}
 \PseudometricError_{\Correspondence_\SequenceIndex}
   (\Pseudometric_{\ComparisonCarrier_\SequenceIndex},\Pseudometric_\ComparisonCarrier)
 \leq2\CoefficientBound\NormComparisonError_\SequenceIndex
       +2\CoordinateComparisonError_\SequenceIndex
          \max_{\substack{\CoefficientVector\in\RealNumbers^\ModelDimension\\\norm{\CoefficientVector}_2=1}}\NormSymbol_\ComparisonCarrier(\CoefficientVector).
\end{equation}
The right-hand side tends to
$0$.
Thus
\begin{equation}\label{eq:local-pseudometric-convergence}
 \PseudometricError_{\Correspondence_\SequenceIndex}
 (\Pseudometric_{\ComparisonCarrier_\SequenceIndex},\Pseudometric_\ComparisonCarrier)\to0.
\end{equation}
If
$\AmbientDisplacement _\SequenceIndex =\sup_{(\ComparisonPoint_\SequenceIndex,\ComparisonPoint)\in\Correspondence_\SequenceIndex}\AmbientMetric (\ComparisonPoint _\SequenceIndex ,\ComparisonPoint )$,
then for every
$(\ComparisonPoint_\SequenceIndex,\ComparisonPoint),
 (\IntegrationPoint_\SequenceIndex,\IntegrationPoint)\in\Correspondence_\SequenceIndex$,
\[
 |\ComparisonMetric _\SequenceIndex(\ComparisonPoint _\SequenceIndex ,\IntegrationPoint _\SequenceIndex )-\ComparisonMetric (\ComparisonPoint ,\IntegrationPoint )|\leq2\AmbientDisplacement _\SequenceIndex .
\]
By \eqref{eq:pseudometric-uniform-comparison}
and \eqref{eq:local-pseudometric-convergence},
we obtain
\[
 |\LocalError \MetricSpace{\ComparisonCarrier _\SequenceIndex }{\ComparisonMetric _\SequenceIndex}-\LocalError \MetricSpace{\ComparisonCarrier }{\ComparisonMetric }|
 \leq\PseudometricError_{\Correspondence_\SequenceIndex}
 (\Pseudometric_{\ComparisonCarrier_\SequenceIndex},\Pseudometric_\ComparisonCarrier)
 +2\AmbientDisplacement _\SequenceIndex \to0.
\]
The bound at
$\MetricSpace{\BaseCarrier }{\MetricSymbol }$
was arranged in Step~\ref{step:local-model-spectral-dimension} by \eqref{eq:finite-spectral-error}.
This proves \ref{item:model-pseudometric}--\ref{item:model-center-error}.
\end{proof}

\paraid{p-ec88a14da3e5}
The construction also represents the isometries of each input space by
orthogonal matrices.

\begin{remark}\label{rem:local-isometry-representation}
\paraid{p-a393c40d9c5a}
Fix a local model from \Cref{thm:local-models} and a representative
$(\FeatureCoordinates_\ComparisonCarrier,\NormSymbol_\ComparisonCarrier)$
constructed from the real orthonormal basis
$\{\BasisFunction_\BasisIndex\}_{0\leq\BasisIndex<\ModelDimension}$
of
$\SpectralSubspace{\ComparisonCarrier}{\SpectralCutoff}$.
Then the construction induces a continuous group homomorphism
\begin{equation}\label{eq:local-isometry-representation-1}
 \IsometryRepresentation{\ComparisonCarrier}{\ComparisonMetric}
 \colon\Isom\MetricSpace{\ComparisonCarrier}{\ComparisonMetric}
       \to\OrthogonalGroup(\ModelDimension),
\end{equation}
with the following matrix entries.
For every
$\FirstIsometry\in\Isom\MetricSpace{\ComparisonCarrier}{\ComparisonMetric}$
and integers
$0\leq\BasisIndex,\SecondBasisIndex<\ModelDimension$,
set
\begin{equation}\label{eq:local-isometry-representation-2}
 \bigl(\IsometryRepresentation{\ComparisonCarrier}{\ComparisonMetric}(\FirstIsometry)\bigr)_{\BasisIndex\SecondBasisIndex}
 =\int_\ComparisonCarrier
   \BasisFunction_\BasisIndex(\FirstIsometry\ComparisonPoint)
   \BasisFunction_\SecondBasisIndex(\ComparisonPoint)
   \,\IntegrationDifferential\SelectedMeasure{\ComparisonCarrier}{\ComparisonMetric}(\ComparisonPoint).
\end{equation}
Indeed,
Step~\ref{step:local-model-isometries} of the proof shows that composition with
$\FirstIsometry$
preserves the spectral subspace and its inner product.
Hence
\begin{equation}\label{eq:isometry-basis-action}
 \BasisFunction_\BasisIndex\circ\FirstIsometry
 =\sum_{\SecondBasisIndex=0}^{\ModelDimension-1}
  \bigl(\IsometryRepresentation{\ComparisonCarrier}{\ComparisonMetric}(\FirstIsometry)\bigr)_{\BasisIndex\SecondBasisIndex}
  \BasisFunction_\SecondBasisIndex.
\end{equation}
For every
$\FirstIsometry,\SecondIsometry\in\Isom\MetricSpace{\ComparisonCarrier}{\ComparisonMetric}$,
applying \eqref{eq:isometry-basis-action} successively to
$\FirstIsometry$
and
$\SecondIsometry$
and using uniqueness of basis coefficients proves
\[
 \IsometryRepresentation{\ComparisonCarrier}{\ComparisonMetric}(\FirstIsometry\circ\SecondIsometry)
 =\IsometryRepresentation{\ComparisonCarrier}{\ComparisonMetric}(\FirstIsometry)
  \IsometryRepresentation{\ComparisonCarrier}{\ComparisonMetric}(\SecondIsometry).
\]
The matrix entries in \eqref{eq:local-isometry-representation-2} are continuous
for the uniform topology on the isometry group.
By \eqref{eq:input-isometry-1},
\eqref{eq:input-isometry-2} and uniqueness of the best approximations,
we obtain the equivariance identity \eqref{eq:local-model-equivariance-1}
and the norm preservation identity \eqref{eq:local-model-equivariance-2}.
For every
$\FirstIsometry\in\Isom\MetricSpace{\ComparisonCarrier}{\ComparisonMetric}$
and
$\ComparisonPoint\in\ComparisonCarrier$,
\begin{equation}\label{eq:local-model-equivariance-1}
 \FeatureCoordinates_\ComparisonCarrier(\FirstIsometry\ComparisonPoint)
 =\IsometryRepresentation{\ComparisonCarrier}{\ComparisonMetric}(\FirstIsometry)
   \FeatureCoordinates_\ComparisonCarrier(\ComparisonPoint).
\end{equation}
For every
$\FirstIsometry\in\Isom\MetricSpace{\ComparisonCarrier}{\ComparisonMetric}$
and
$\CoefficientVector\in\RealNumbers^\ModelDimension$,
\begin{equation}\label{eq:local-model-equivariance-2}
 \NormSymbol_\ComparisonCarrier
   \bigl(\IsometryRepresentation{\ComparisonCarrier}{\ComparisonMetric}(\FirstIsometry)\CoefficientVector\bigr)
 =\NormSymbol_\ComparisonCarrier(\CoefficientVector).
\end{equation}
A change of basis by
$\OrthogonalChange\in\OrthogonalGroup(\ModelDimension)$
replaces this representation by
$\FirstIsometry\mapsto\OrthogonalChange
 \IsometryRepresentation{\ComparisonCarrier}{\ComparisonMetric}(\FirstIsometry)
 \OrthogonalChange^{-1}$.
For the constant model used at a singleton center,
take the trivial homomorphism into
$\OrthogonalGroup(1)$.
For a compact group
$\ActingGroup$,
the Peter--Weyl theorem asserts that the matrix coefficients of finite-dimensional
continuous unitary representations have dense linear span in
$\ContinuousFunctions(\ActingGroup,\ComplexNumbers)$
\cite[Chapter~6, Section~2, pp.~165--167]{DiestelSpalsbury2014}.
In that proof,
convolution operators commute with left translations,
and their nonzero eigenspaces carry finite-dimensional unitary representations.
Here the distance operator commutes with the action of
$\Isom\MetricSpace{\ComparisonCarrier}{\ComparisonMetric}$,
so its invariant spectral subspace defines
$\IsometryRepresentation{\ComparisonCarrier}{\ComparisonMetric}$.
\Cref{lem:spectral-continuity} permits comparison of these subspaces as the compact
metric space varies.
\end{remark}

\begin{remark}\label{rem:mds-comparison}
\paraid{p-1384428ed1f6}
For compact metric measure spaces,
Kroshnin,
Stepanov,
and Trevisan
\cite[arXiv v2, Theorem~5.4 and Remark~5.5]{KroshninStepanovTrevisan2022}
prove stability of multidimensional scaling coordinates in
$\LebesgueSpace^2$
using optimal couplings for the Gromov--Wasserstein distance of order four
and minimizing over orthogonal transformations,
assuming convergence for that distance and a gap
after the last retained positive eigenvalue.
\Cref{thm:local-models} uses best
$\LebesgueSpace^\IntegrabilityExponent$
approximations of distance functions and varying norms to approximate
the metric uniformly.
\end{remark}
